\documentclass[../../../include/open-logic-section]{subfiles}

\begin{document}

\olfileid[pt]{sth}{replacement}{intro}
\olsection{Introdução}

A Substituição é o esquema axiomático que distingue $\ZF$
de~$\Z$. Recorremos a ele ao longo de
\crefrange{sth:ordinals::chap}{sth:spine::chap}. Neste capítulo,
abordaremos finalmente a questão: a Substituição está justificada? 

Para tornar a questão precisa, convém observar que a Substituição é
realmente bastante \emph{forte}. Teremos uma ideia de quão forte é ao longo deste capítulo (e de novo na \olref[sth][card-arithmetic][fix]{sec}). Isto sugere que uma justificação é de fato necessária. 

Discutiremos dois tipos de justificação. Em traços largos: uma justificação \emph{extrínseca} é uma tentativa de justificar um axioma pelos seus frutos; uma justificação \emph{intrínseca} é uma tentativa de justificar um axioma sugerindo que ele é vindicado pelos conceitos matemáticos em causa. Teremos uma ideia mais clara do que isto significa ao longo deste capítulo, mas não passa da ponta de um icebergue. Para mais pormenores, veja em particular \citeauthor{Maddy1988a} (\citeyear{Maddy1988a} e \citeyear{Maddy1988b}).

\end{document}
